\subsection{The arithmetic intersections for the three CM fibres}

\begin{statement}{AF.1}{Proved}{A quadratic-by-quartic intersection bound}
Let $K/\Q$ be quadratic and $L/\Q$ cyclic quartic, with $K\cap L=\Q$.
Inside a fixed algebraic closure, put $M=K^{\mathrm{ab}}\cap L^{\mathrm{ab}}$.
Then $M/\Q^{\mathrm{ab}}$ has degree at most two.
\end{statement}
\begin{proof}
Both maximal abelian extensions are Galois over $\Q$, and contain
$\Q^{\mathrm{ab}}$. In $G=\Gal(M/\Q)$ put
$A=\Gal(M/K)$ and $B=\Gal(M/L)$. These are closed abelian normal
subgroups of indices two and four. Since $K\cap L=\Q$, we have
$AB=G$. The subgroup $Z=A\cap B$ is central: its elements commute
with $A$ and $B$, and these generate $G$. The quotient $G/Z$ is
$\Gal(KL/\Q)=C_2\times C_4$. Choose $a\in A$ generating its
order-four factor and $b\in B$ generating its order-two factor.
Then $G$ is topologically generated by $a,b,Z$, and all commutators
are central and generated by $[a,b]$. Because $b^2\in Z$,
$[a,b]^2=[a,b^2]=1$. Thus $G'$ has order at most two. The maximal
abelian subextension of $M/\Q$ is exactly $\Q^{\mathrm{ab}}$,
since that field is already contained in $M$. This proves the bound,
including for the infinite profinite extensions.
\end{proof}

\begin{statement}{AF.2}{Proved}{The Gaussian--quintic intersection}
For $K_4=\Q(i)$ and $K_5=\Q(\zeta_5)$,
\[
 K_4^{\mathrm{ab}}\cap K_5^{\mathrm{ab}}
       =\Q^{\mathrm{ab}}(\sqrt[4]{5}).
\]
\end{statement}
\begin{proof}
Let $D_5=\Q(i,\sqrt[4]{5})$. The polynomial $X^4-5$ is Eisenstein
at five, and its real root generates a nonnormal real quartic field.
Thus $D_5/\Q$ has dihedral group of order eight. Its three quadratic
subfields are $\Q(i)$, $\Q(\sqrt5)$ and $\Q(\sqrt{-5})$.
Its extension over $K_4$ is cyclic quartic. The only quadratic
subfield of $K_5$ is $\Q(\sqrt5)$, and $K_5$ cannot be a quartic
Galois subfield of $D_5$, because the abelianization of the dihedral
group has exponent two. Consequently
$D_5\cap K_5=\Q(\sqrt5)$; the group over this intersection is the
Klein four group. Hence $D_5K_5/K_5$ is abelian too.
It follows that $\Q^{\mathrm{ab}}(\sqrt[4]5)$ lies in the claimed
intersection. It is quadratic over $\Q^{\mathrm{ab}}$: its square
is already there, whereas the fourth root itself is not, since
$D_5/\Q$ is nonabelian. AF.1 supplies the opposite bound.
\end{proof}

\begin{statement}{AF.3}{Proved}{The Eisenstein--quintic intersection and a finite defining field}
Let $K_3=\Q(\sqrt{-3})$ and $E_3:y^2=x^3+1$. The group of
$D_{35}=\Q(E_3[5])$ over $\Q$ is $C_{24}\rtimes_5 C_2$.
Let $M_{35}$ be its subfield fixed by the unique subgroup of
order three in the Cartan $C_{24}$. Then
\[
 \Gal(M_{35}/\Q)=C_8\rtimes_5 C_2,\qquad
 K_3^{\mathrm{ab}}\cap K_5^{\mathrm{ab}}
                    =\Q^{\mathrm{ab}}M_{35}.
\]
The extension on the right has degree two over $\Q^{\mathrm{ab}}$.
This description specifies $M_{35}$ unambiguously as a finite
subfield of the division field, and does not depend on a choice
of its primitive generator.
\end{statement}
\begin{proof}
The endomorphism $(x,y)\mapsto(\zeta_3x,y)$ is defined over $K_3$.
The $\mathbf F_5$-space $E_3[5]$ is one-dimensional over
$\mathcal O_{K_3}/5=\mathbf F_{25}$. Therefore the image of
$G_{K_3}$ is contained in its cyclic unit group of order 24.
The field $D_{35}$ contains $K_3$. Indeed, if a Galois element
trivial on $E_3[5]$ conjugated this endomorphism to its conjugate,
their difference would annihilate $E_3[5]$. That difference has
CM norm three, so it is invertible on five-torsion, a contradiction.

We now prove the full Cartan image, rather than assume a general
surjectivity theorem. Counting gives $\#E_3(\mathbf F_7)=12$.
Frobenius at either prime of $K_3$ above seven has trace $-4$
and norm seven, so it is multiplication by $-2+\sqrt{-3}$ or its
conjugate. In $\mathbf F_{25}=\mathbf F_5[t]/(t^2-2)$ these two
units $(3,1)$ and $(3,-1)$ have order 24. The attached exact
program checks the point count and both orders. Hence the image
is all of $C_{24}$. Complex conjugation has order two, acts on
this Cartan by the nontrivial field automorphism, and thus acts
by the exponent five. This proves the asserted order-48 group.
For a polynomial description of the finite division field, its
$x$-coordinates are the roots of
\[
       5X^{12}+380X^9-240X^6-1600X^3-256;
\]
one also adjoins their $y$-coordinates satisfying $y^2=x^3+1$.

Write the order-16 quotient as
$\langle b,c\mid b^8=c^2=1,\ cbc^{-1}=b^5\rangle$.
Its derived subgroup is $\langle b^4\rangle$ of order two.
The extension over $K_3$ is cyclic, generated by $b$.
The Weil pairing puts $K_5$ in $D_{35}$, and the odd subgroup
of order three fixes it. Choose the generator so that its
determinant is the chosen generator of $\mathbf F_5^\times$;
complex conjugation has determinant $-1$. The determinant kernel
in the order-16 group is
\[
       \{1,b^4,b^2c,b^6c\}=\langle b^2c\rangle\simeq C_4.
\]
Thus $M_{35}$ is abelian over both $K_3$ and $K_5$, but not over
$\Q$. Its compositum with $\Q^{\mathrm{ab}}$ is a nontrivial
quadratic extension, because its commutator subgroup has order
two. AF.1 proves equality with the entire intersection.
\end{proof}

\begin{statement}{AF.4}{Known}{The full elliptic torsion fields and the CM upper bound for the Jacobian}
Let $J_5=\operatorname{Jac}(y^2=x^5-1)$ and define Galois extensions
\[
 L_3=\Q^{\mathrm{ab}}(E_3[\mathrm{tors}]),\quad
 L_4=\Q^{\mathrm{ab}}(E_4[\mathrm{tors}]),\quad
 L_5=\Q^{\mathrm{ab}}(J_5[\mathrm{tors}]),
 \qquad E_4:y^2=x^3-x.
\]
Then $L_3=K_3^{\mathrm{ab}}$, $L_4=K_4^{\mathrm{ab}}$, and
$L_5\subseteq K_5^{\mathrm{ab}}$.
\end{statement}
\begin{proof}[Derivation from the stated CM theorem]
We use Weber's ray-class-field theorem in the precise form of
\citet[Theorem 8]{ClarkCookStankewicz2009}: for an elliptic curve
with the maximal CM order, adjoining its $j$-invariant and the
values of a Weber function on its $N$-torsion gives the ray class
field of modulus $N$. Here the Weber functions can be taken to
be $x^3$ for $E_3$ and $x^2$ for $E_4$. Both CM fields have class
number one, and the displayed curves and all their endomorphisms
are defined over those fields. Their full torsion fields are
abelian over the CM field: each Tate module is a rank-one module
over the corresponding completed CM order and Galois acts by
units. Every ray class field therefore lies in the torsion field,
and every finite abelian extension lies in a ray class field.
This proves the two equalities.

For the genus-two Jacobian the automorphism
$(x,y)\mapsto(\zeta_5x,y)$ induces an action of
$\mathcal O_{K_5}=\Z[\zeta_5]$ defined over $K_5$. On its
four-dimensional rational first homology this action has minimal
polynomial $\Phi_5$: the holomorphic differentials $dx/y$ and
$x\,dx/y$ have eigenvalues $\zeta_5,\zeta_5^2$, and their
conjugates supply the other two. Each Tate module is consequently
rank one over $\mathcal O_{K_5}\otimes\Z_\ell$. More explicitly,
a stable full lattice is torsion-free of rank one over each
local discrete valuation factor, hence free there, including at
$\ell=5$. Galois over $K_5$ acts by units of this commutative
ring. Every torsion extension over $K_5$ is abelian. Since
$\Q^{\mathrm{ab}}\subset K_5^{\mathrm{ab}}$, the asserted
containment follows. No equality $L_5=K_5^{\mathrm{ab}}$ is used.
\end{proof}

\begin{statement}{AF.5}{Proved}{The Jacobian contains both predicted extra quadratic intersections}
The field $L_5$ contains $\sqrt[4]5$ and $\Q^{\mathrm{ab}}M_{35}$.
\end{statement}
\begin{proof}
For the first assertion we use the specialization statement,
not a generic-curve assertion: \citet[Corollary 1.2(a)]{Yelton2019}
places in the eight-division field the fourth roots of
$\prod_{j\ne i}(\alpha_j-\alpha_i)$, after the stated cyclotomic
adjoining, for the distinct roots $\alpha_i$ of an odd-degree
hyperelliptic polynomial. The necessary cyclotomic constants
are in $\Q^{\mathrm{ab}}$. At the root $\alpha_i=1$ of
$x^5-1$ this product is $f'(1)=5$. Thus
$\sqrt[4]5\in L_5$.

For the second assertion, put $D_J=\Q(J_5[3])$.
It contains $K_5$: a Galois element fixing the three-torsion
cannot change the endomorphism $\zeta_5$ to $\zeta_5^a$ for
$a\ne1$, since their difference has norm five and is invertible
on three-torsion. Let $U=\Gal(D_J/K_5)$, viewed through the
faithful rank-one representation as a subgroup of
$(\mathcal O_{K_5}/3)^\times=\mathbf F_{81}^\times$.
The Rosati involution for the canonical principal polarization
is complex conjugation, which is $u\mapsto u^9$ on this field.
The Weil pairing therefore gives
\[
                 u^{10}=u\overline u=\chi_3(u)\in\mathbf F_3^\times.
\]
So $U$ lies in the cyclic group of order 20 cut out by $u^{20}=1$.
The prime eleven splits in $K_5$, and is a prime of good reduction
of this curve. Frobenius at a prime above eleven belongs to
$G_{K_5}$ and has multiplier $11=-1\pmod3$. Its image satisfies
$u^{10}=-1$, forcing its two-primary part to have order four.
Thus the two-primary part of $U$ is exactly $C_4$.

Let $O$ be the unique odd-order subgroup of $U$. It is
characteristic in $U$, hence normal in $\Gal(D_J/\Q)$; it has
order one or five. Set $M_J=D_J^O$. There is an exact sequence
\[
  1\longrightarrow\langle u\rangle\simeq C_4
   \longrightarrow\Gal(M_J/\Q)\longrightarrow C_4
   \longrightarrow1.
\]
A generator of $\Gal(K_5/\Q)$ sends $\zeta_5$ to $\zeta_5^2$.
On $\mathbf F_{81}$ this is the third power of Frobenius, since
$3^3=2\pmod5$; on its subgroup $C_4$ it acts by exponent
$27=-1\pmod4$. Consequently the displayed group is nonabelian,
its derived subgroup is exactly $\langle u^2\rangle$, and that
subgroup is central. These conclusions do not require the exact
sequence to split.

The Weil pairing puts $\mu_3$ in $D_J$, and the odd group $O$
fixes it, so $K_3\subset M_J$. The multiplier character on
$\langle u\rangle$ sends its generator to $-1$. Choose a lift
$v$ of a generator of the quotient $C_4$ and, if necessary,
replace it by $uv$ to make its multiplier one. The multiplier
kernel is generated by $u^2$ and $v$: it has index two, since
these generate four quotient cosets and the central order-two
subgroup. The two generators commute, so this kernel is abelian.
It is exactly $\Gal(M_J/K_3)$. The group over $K_5$ is already
cyclic. Thus $M_J$ lies in $K_3^{\mathrm{ab}}\cap K_5^{\mathrm{ab}}$,
and its compositum with $\Q^{\mathrm{ab}}$ is the nontrivial
quadratic extension described in AF.3. Since $M_J\subset D_J$,
that whole compositum lies in $L_5$.

The finite-field computations, both possible order-16 extension
presentations, their commutator and multiplier kernels, and the
Frobenius certificate in AF.3 are checked in
\texttt{code/check\_cm\_fields.py}. They support the exact group
arguments; they do not substitute for the CM and specialization
comparisons proved above.
\end{proof}

\begin{statement}{AF.6}{Proved}{All three pairwise torsion-field intersections and the triple intersection}
With the definitions in AF.4, the complete arithmetic answer is
\[
\begin{aligned}
 L_3\cap L_4&=\Q^{\mathrm{ab}}(\sqrt[4]3),\\
 L_4\cap L_5&=\Q^{\mathrm{ab}}(\sqrt[4]5),\\
 L_3\cap L_5&=\Q^{\mathrm{ab}}M_{35},\\
 L_3\cap L_4\cap L_5&=\Q^{\mathrm{ab}}.
\end{aligned}
\]
The three nontrivial quadratic extensions in this display are
linearly disjoint over $\Q^{\mathrm{ab}}$ and have compositum
of degree eight. These are arithmetic field statements, not a
construction of three characters of $\GT$.
\end{statement}
\begin{proof}
The first equality is A.12; its finite radical may be written
$\sqrt[4]3$ because $\sqrt2\in\Q^{\mathrm{ab}}$.
The second and third follow from AF.2--AF.5: the maximal
abelian intersections give upper bounds, and AF.5 supplies the
full quadratic subextensions inside the actual Jacobian field.

For the triple intersection it suffices to show that
$\sqrt[4]3\notin K_5^{\mathrm{ab}}$. Let
$D_3=\Q(i,\sqrt[4]3)$, whose dihedral Galois group has order
eight and whose quadratic subfields are $\Q(i)$,
$\Q(\sqrt3)$ and $\Q(\sqrt{-3})$.
Its intersection with $K_5$ is $\Q$: it contains no
$\Q(\sqrt5)$ and has no cyclic quartic quotient.
Thus $D_3K_5/K_5$ is still nonabelian. If the fourth root
belonged to $K_5^{\mathrm{ab}}$, then so would $i$, all of
$D_3$, and their compositum with $K_5$, a contradiction.
The triple intersection is therefore trivial over
$\Q^{\mathrm{ab}}$.

Finally the first and second quadratic fields are distinct,
since the first is not contained in $K_5^{\mathrm{ab}}$ whereas
the second is. Both lie in $L_4$, so their biquadratic
compositum lies in $L_4$. The third does not lie in $L_4$,
by the triple-intersection assertion, because it already lies
in $L_3\cap L_5$. It is therefore independent of the first
two, proving the degree-eight assertion.
\end{proof}

\begin{statement}{AF.7}{Open}{The missing passage from these fields to order-five frame characters}
The equalities in AF.6 give relations among the actual Galois
representations of the three fibres. They do not provide a
corrected homomorphism from the multiplier-one subgroup of
$\GT$ to an order-five norm-one target, or a comparison of that
homomorphism with the Jacobian representation. We have not
proved either of those maps. Consequently we have not extended
the three arithmetic quadratic relations to characters of
$\GT$, nor evaluated such proposed characters away from Galois.
\end{statement}
The smallest missing comparison is already the three-torsion
coordinate of $J_5$ used in AF.5: we would need a genuine
five-point action, a specified path and normalization of its
order-five frame, and a proof that its finite coordinate has
the arithmetic quotient $M_J$ above. The usual genus-two braid
representation fails the sphere relation (A.17), so it cannot
be used without a different group and descent argument.
